\documentclass[11pt]{article}
\usepackage[margin=1in]{geometry}
\usepackage{natbib}
\usepackage{hyperref}
\usepackage{booktabs}

\title{Experimental Plans}
\author{}
\date{}

\begin{document}
\maketitle

% Methods reference for this project. Records final design decisions and their
% rationale, not a change log -- see git history for how decisions evolved.

\section{Model and Training Pipeline}

\subsection{Architecture and Tokenizer}

\begin{itemize}
  \item Model: GPT-2-small, 12 layers, 768 hidden dim, 12 attention heads,
    $\sim$117M parameters \citep{radford2019language}. GPT-2-medium (24 layers,
    1024 hidden dim, 16 heads, $\sim$345M) is reserved for confirmatory runs on
    conditions selected by the pilot in \S\ref{sec:seeds}.
  \item Tokenizer: pretrained GPT-2 BPE, used unmodified.
  \item Context length: 512 tokens. Positional embeddings are trained from
    scratch, so context length is not constrained by the pretrained GPT-2
    checkpoint.
\end{itemize}

\subsection{Corpus}
\label{sec:corpus}

\begin{itemize}
  \item \textbf{Backbone: OpenWebText} \citep{Gokaslan2019OpenWebText},
    $\approx$9.04B tokens, comfortably large enough for a single non-repeating
    pass at the token budget below.
  \item \textbf{Streamed and consumed in bounded chunks, never downloaded in
    full}: a storage-constraint requirement, not a performance optimization.
    The backbone loader (\texttt{preprocess/backbone\_openwebtext.py}) pulls
    documents from the \texttt{Skylion007/openwebtext} HF dataset via
    \texttt{datasets.load\_dataset(..., streaming=True)}, buffers raw text in
    memory only until a chunk boundary (default 2GB, configurable) is reached,
    hands that chunk to the assembler for tokenizing and packing
    (\S\ref{sec:assembly}), then discards the buffer -- no more than one
    chunk's raw text is held at once, and no persistent local copy of the
    dataset is written to disk.
  \item Injected data: synthetic biography documents, one per entity
    (\S\ref{sec:entity}), inserted at randomly scheduled positions in the
    backbone token stream (\S\ref{sec:scheduler}).
  \item \textbf{Independent variable: exposure count only.} Each entity has
    exactly one contested relation with two candidate values $A$ and $B$; the
    swept quantity is the split $(n_A, n_B)$ of a fixed total exposure budget
    $T$ between the two sides (\S\ref{sec:scheduler}, \S\ref{sec:scale}).
    Recency is not tracked or manipulated: in a from-scratch training run, an
    occurrence's position late in the stream is confounded with the cosine LR
    schedule's decay toward that point, so treating recency as a variable
    would entangle exposure timing with optimization dynamics.
  \item \textbf{Token budget: fixed at $2.5\times10^9$ total training tokens},
    the Chinchilla-optimal point for a 124M-parameter model
    \citep{hoffmann2022chinchilla} at $\approx$20 tokens/parameter. Per
    \S\ref{sec:corpus-vs-training-tokens}'s distinction, this is the total
    training-token figure (batch $\times$ context $\times$ steps,
    \S\ref{sec:optimization}), not necessarily the on-disk corpus size; since
    OpenWebText has enough unique tokens to cover $2.5\times10^9$ in a single
    pass with no repetition, \texttt{total\_tokens} (on-disk, \S\ref{sec:assembly})
    is set equal to the training-token target, so the run needs zero epochs
    of repetition over its own backbone -- the cleanest point on
    \citet{muennighoff2023scaling}'s scaling curve.
\end{itemize}

\subsection{Entity and Document Construction}
\label{sec:entity}

\begin{itemize}
  \item \textbf{Entity schema}: each synthetic entity is a tuple
    $(\text{name}, \text{contested}, \{\text{background}_j\})$. \emph{Name}
    is drawn from a globally unique entity-name pool, disjoint from every
    other human-name-shaped value pool. \emph{Contested} is one relation type
    assigned two candidate values $\text{val}_A, \text{val}_B$ from that
    relation's value pool and a split $(n_A, n_B)$, $n_A + n_B = T$
    (\S\ref{sec:scheduler}). $\{\text{background}_j\}$ is a fixed-size set of
    $K=4$ further relation-type/value pairs, each assigned a single value
    that never varies. Together these give each biography 5 attributes
    (1 contested $+$ 4 background), paralleling Allen-Zhu \& Li's bioS
    structure. Background content is restated every time \emph{any}
    occurrence of the entity's biography renders (\S\ref{sec:scheduler}), so
    its exposure count is exactly $T$, not independently controlled.
  \item \textbf{Relation-type assignment, balanced by construction}: one
    shuffled list of (relation type, split level) pairs is built, each of the
    14 $\times$ (number of split levels) combinations repeated the same
    replicate count (\S\ref{sec:scale}); one pair is assigned per entity by
    walking the shuffled list, guaranteeing every relation type is used as
    the contested type equally often, both overall and within every split
    level. Background types follow a fixed cyclic offset from the contested
    type: given a fixed ordering of the 14 relation types, an entity whose
    contested type sits at position $c$ takes background types at positions
    $c{+}1, c{+}2, c{+}3, c{+}4$ (mod 14), which balances background-type
    usage without separate randomization. Values are drawn by walking a
    shuffled copy of the relevant pool in order, reshuffling and starting a
    new lap once exhausted, so usage counts differ by at most one; the
    contested pair's $\text{val}_A, \text{val}_B$ are drawn this way and
    checked for surface-form divergence (\S\ref{sec:divergence}) before being
    accepted.
  \item \textbf{One contested relation per entity, not several}: whenever an
    entity's biography renders to advance one contested relation's exposure
    count, every other attribute in that instance is restated too. Allowing
    multiple simultaneously-contested relations per entity would couple
    their exposure patterns through shared document instances, correlating
    which side wins across relations by construction. A single contested
    relation per entity keeps each conflict pair's exposure pattern clean.
  \item \textbf{Relation-type inventory} (14 types, spanning affiliation,
    fixed life-event, and ongoing-state semantics so findings are not an
    artifact of one relation's surface statistics): alma mater; employer;
    birthplace; current residence (distinct value space from birthplace);
    award or honor received; authored work or invention; professional or
    hobbyist affiliation; field of expertise; mentor or doctoral advisor
    (named individual, drawn from a name pool disjoint from all entity
    names); primary funding source; professional license or certification;
    publication venue; civic or volunteer organization; working or
    professional language. Each relation type has its own value pool,
    checked against named entities in WikiText-103 and against every other
    value pool. \textbf{Empirically, only 7 of these 14 reach reliable
    recall} at the trained checkpoint; intervention experiments restrict the
    contested-relation-type inventory to that subset
    (\S\ref{sec:relation-restriction}).
  \item \textbf{Value-pool reuse policy}: a value may be reused across
    entities if it is not itself shaped like a human name (institution,
    place, award, license, venue values are drawn from finite pools and may
    repeat, matching real-corpus statistics). Any value that is itself a
    person's name -- the entity's own name, or the mentor/advisor relation's
    value -- is drawn from a globally unique pool and never reused. The
    entity-name pool and every other human-name-shaped value pool are each a
    separate bank, checked for mutual disjointness.
  \item \textbf{Document rendering}: LLM-authored, not mechanically assembled
    from templates. Operationalized as bioS(N) in spirit
    \citep{allenzhu2023physics} -- a full biography per (entity, side) -- but
    the text itself is written by an auxiliary LLM (\S\ref{sec:generation})
    given the entity's name and its 5 facts (contested relation resolved to
    one side, plus all $K$ background relations) handed over \emph{flat}: the
    prompt gives no indication that one fact is contested, so the model has
    no opportunity to hedge or write differently about it than about an
    ordinary background fact. One call per (entity, side) produces $V$
    independent variants of a short, coherent paragraph stating all 5 facts
    (\S\ref{sec:generation}). There is no template bank per relation type and
    no first-mention/restatement split; the model composes the whole
    paragraph, choosing its own structure and word order freely each variant.
  \item \textbf{Value fidelity is the one hard constraint on that freedom}:
    every fact's value must appear verbatim in the generated text (enforced
    by prompt instruction and checked post-hoc, \S\ref{sec:generation}),
    because per-fact recall (\S\ref{sec:eval}) is scored by locating that
    exact string and reading the model's probability at that point. The
    post-hoc check is case-\emph{insensitive}: real common-noun values
    (e.g.\ a \texttt{field\_expertise} term) are reliably case-folded
    mid-sentence by a fluent writer, unlike invented proper-noun values,
    which keep their capitalization; exact case was never the actual
    requirement, only that the value's position in the text is findable for
    scoring.
  \item \textbf{Values must be name-shaped, not full descriptive phrases}: a
    fluent LLM writer reliably keeps short, name-shaped values verbatim (a
    place, an award, a society name) but naturally rewords a full descriptive
    clause into different prose (e.g.\ ``a senior analyst at Ashgrove
    Dynamics'' becomes ``works at Ashgrove Dynamics as a senior analyst''),
    silently breaking value fidelity. \texttt{alma\_mater\_values},
    \texttt{employer\_role\_values}, and \texttt{civic\_role\_values} are
    accordingly shortened to bare institution/organization names. Any future
    value pool must be name-shaped for the same reason, not merely
    fictional.
  \item \textbf{Variant choice is round-robin, not resampled}: at each
    occurrence event (\S\ref{sec:scheduler}), one of that side's
    pre-generated variants is used, drawn without repetition until all have
    been used once, then reshuffled, so an entity's occurrences do not reuse
    the same paragraph back-to-back purely by chance. State is kept per
    entity, persisting across that entity's occurrence events.
  \item \label{sec:variants-per-side} \textbf{Total variants per entity fixed
    at 8, allocated unevenly across sides} via \texttt{variants\_per\_side}
    (\texttt{preprocess/scheduler.py}). Given the entity's own $(n_A, n_B)$
    (\S\ref{sec:scale}), $v_B = \max(1, \mathrm{round}(8 \cdot n_B /
    (n_A{+}n_B)))$ and $v_A = 8 - v_B$: the side with more scheduled
    occurrences gets proportionally more fresh material before round-robin
    repeats kick in. At the balanced split $(40,40)$ this gives $(4,4)$; at
    every split from $(68,12)$ onward the low side floors to a single
    variant. Combined with $T=80$ (\S\ref{sec:scale}), most occurrences
    necessarily reuse a small number of paragraphs many times over -- an
    explicit, accepted design choice (\S\ref{sec:assembly}).
  \item \textbf{Recall scoring is decoupled from training text}: side $A$
    and side $B$ training text are independently LLM-authored free text with
    no common stem, so \S\ref{sec:divergence}'s single-token surface-form
    divergence guarantee (construction-guaranteed under shared templates)
    does not carry over. Per-fact recall (\S\ref{sec:eval}) is instead scored
    on a separate, fixed template-based probe, not on training vignettes --
    \S\ref{sec:assembly} has the full resolution.
\end{itemize}

\subsection{Pool and Vignette Generation}
\label{sec:generation}

Every name/value pool referenced in \S\ref{sec:entity}, and every entity's
biography text itself, is produced by an auxiliary LLM (gpt-4o-mini via the
OpenAI API) rather than hand-authored, so pool/vignette content does not
itself encode the experimenters' assumptions about frequency. Pool generation
uses one meta-prompt per pool (15: the entity-name and mentor-name pools, and
thirteen relation-type value pools; Table~\ref{tab:pools}), substituting only
a category description and a target count into an otherwise fixed template:

\begin{quote}\small\ttfamily
Generate \{n\} distinct entries for this category: \{description\}\newline
Rules: each entry is \{name-shaped / short-phrase\}; \{fictional-required /
real-terms-allowed\}; no trailing punctuation/numbering/commentary; prefer a
distinctive first word per entry. [Do not repeat: \{already-used entries\}.]\newline
Output ONLY a JSON array of exactly \{n\} strings.
\end{quote}

Vignette generation uses one call per (entity, side): the entity's name and
its 5 facts (contested relation resolved to that side, plus all $K$
background relations) are substituted flat into a fixed template, with no
indication that any one fact is contested:

\begin{quote}\small\ttfamily
Write a short biography of a fictional person for a synthetic training corpus.\newline
Person: \{name\}\newline
Facts (state every one of these somewhere in the biography):\newline
- \{relation label\}: \{value\} [$\times$5, one line per fact]\newline
Write \{n\_variants\} independent variants of this biography.\newline
Rules: each variant is a short, coherent paragraph (3--6 sentences), third person,
plain declarative prose; mentions \{name\}'s full name at least once and states every
fact above; each fact's value must appear verbatim, exactly as given (automatic
scoring later depends on finding the exact string); presents every fact flatly, as
simply true, with no hedging or editorializing; the \{n\_variants\} variants are
genuinely different from each other (opening fact, sentence order, phrasing), not
synonym swaps of one paragraph; invents no additional facts; fictional/generic
otherwise.\newline
Output ONLY a JSON array of exactly \{n\_variants\} strings, one full paragraph per
variant.
\end{quote}

\textbf{Generation procedure (pools)}: requests are issued in batches, each
sized at $\lceil 1.2 \times \text{remaining}\rceil$ (20\% leeway to absorb
rejected duplicate/invalid entries without a guaranteed extra round-trip) and
capped at 150 entries per call; failed calls retry up to 5 times with
exponential backoff. A post-generation cleaning pass removes exact
duplicates, malformed entries, and, for name pools, entries not shaped like
``Firstname Lastname''; the cross-pool and WikiText-103 disjointness checks of
\S\ref{sec:entity} run afterward, on the cleaned pools.

\textbf{Generation procedure (vignettes)}: one call per (entity, side),
requesting that side's target variant count
(\texttt{v\_a}/\texttt{v\_b} from \S\ref{sec:entity}'s
\texttt{variants\_per\_side}), sampled at a moderately elevated temperature
(higher than the near-deterministic pool/template calls) so the $V$ variants
are genuinely structurally distinct. Each returned variant is checked
post-hoc for value fidelity (exact substring match for every one of the 5
values); a failing variant is dropped and the call retried up to a bounded
number of attempts. Value pools must be name-shaped (not full descriptive
phrases) for this check to hold reliably -- a lesson learned from a
population-scale generation pass, where one full-phrase pool
(\texttt{civic\_role\_values}, e.g.\ ``board member of X'') was the sole
cause of a small cluster of vignette-generation failures, fixed by shortening
that pool to bare organization names.

\textbf{Concurrency}: at population scale ($\sim$1{,}680 entities $\times$ 2
sides $=$ 3{,}360 calls) requests are issued concurrently (async, bounded by
a semaphore, default cap 10 in-flight), with per-request exponential backoff
for transient/rate-limit errors. The run is resumable per entity -- an
(entity\_id, side) pair already written to disk is skipped -- per CLAUDE.md
\S5's restart-safety requirement for long-running jobs.

\begin{table}[h]
\centering\small
\begin{tabular}{lrl}
\toprule
Pool & Target size & Notes \\
\midrule
entity\_names & 2000 & globally unique \\
mentor\_names & 2000 & globally unique, disjoint from entity\_names \\
alma\_mater\_values & 150 & reusable across entities \\
employer\_role\_values & 150 & reusable across entities \\
birthplace\_values & 150 & reusable; disjoint from current\_residence\_values \\
current\_residence\_values & 150 & reusable; disjoint from birthplace\_values \\
award\_honor\_values & 150 & reusable across entities \\
authored\_work\_values & 150 & reusable across entities \\
affiliation\_values & 150 & reusable across entities \\
field\_expertise\_values & 150 & real terms; no WikiText-103 overlap check \\
funding\_source\_values & 150 & reusable across entities \\
license\_certification\_values & 150 & reusable across entities \\
publication\_venue\_values & 150 & reusable across entities \\
civic\_role\_values & 150 & reusable across entities \\
working\_language\_values & 60 & real terms; no WikiText-103 overlap check \\
\bottomrule
\end{tabular}
\caption{Generated pools (\texttt{preprocess/schema.py}). Each entity draws 5 of
these 14 relation types at random (\S\ref{sec:entity}), so population size
(\S\ref{sec:scale}) is set independently of pool size.}
\label{tab:pools}
\end{table}

\subsection{Surface-Form Divergence}
\label{sec:divergence}

Per-fact recall (\S\ref{sec:eval}) is scored at a single token to avoid a
length confound: comparing joint log-probability across candidate values of
differing token count penalizes the longer candidate independent of any
frequency effect. Scoring at the first token where two candidate values
diverge avoids this, provided two requirements are enforced at generation
time:

\begin{itemize}
  \item \textbf{Construction}: templates are authored so no shared filler
    token sits between the template stem and the value -- the token stream
    diverges at the value itself.
  \item \textbf{Verification}: checked per rendered context, not per bare
    string, since BPE assigns a different token to the same word depending
    on leading-space / sentence-initial position.
\end{itemize}

Verification is three on-demand checks, applied to the single contested pair
assigned to an entity, rather than an $O(n^2)$ pairwise pass over the whole
value pool (an earlier, stricter version of this check rejected an
unworkable fraction of large pools):
\begin{itemize}
  \item[(a)] \textbf{On-demand pairwise divergence}, applied to the two
    values assigned as an entity's contested pair, at assignment time.
  \item[(b)] \textbf{Exact-duplicate ban}, enforced unconditionally at
    generation time (\S\ref{sec:generation}).
  \item[(c)] \textbf{No overly common first token}: a value is rejected if
    its distinguishing first token (post-stem) decodes to fewer than 3
    characters. Rare/invented words disproportionately fall back to such
    short BPE subwords, which are shared by a large fraction of the
    vocabulary and therefore likely to collide with whatever a value is
    eventually paired against.
\end{itemize}

These checks gate which val$_A$/val$_B$ pairs are offered to the vignette
generator (\S\ref{sec:generation}) and hold for the shared eval-time
\emph{template} (\S\ref{sec:assembly}). They do not guarantee that
independently-authored \emph{training} text diverges at one shared token
position, since side $A$'s and side $B$'s paragraphs may differ in wording
well before the value; per-fact recall is accordingly scored on the separate
eval-time probing format, not on training text (\S\ref{sec:assembly}).

\subsection{Exposure Scheduler}
\label{sec:scheduler}

Fact placement is generated by a scheduler decoupled from document content,
mapping each entity's contested-relation split $(n_A, n_B)$ to token
positions in the full backbone-plus-injection stream:

\begin{itemize}
  \item \textbf{Two threads, one document}: side $A$ and side $B$ of an
    entity's contested relation are two \emph{threads} that share the same
    underlying biography-rendering machinery (\S\ref{sec:entity}), rather
    than two independent documents. Thread $A$ contributes $n_A$ occurrence
    events, thread $B$ contributes $n_B$, $n_A + n_B = T$.
  \item \textbf{Placement}: all $T$ occurrence events for an entity are
    placed i.i.d.\ uniformly at random across the full planned
    training-token stream, then labeled $A$ or $B$ by thread. Placement
    timing is not a studied variable (\S\ref{sec:corpus}), so a single
    well-understood uniform randomization is used rather than any
    position-dependent mechanism.
  \item \textbf{Rendering coupling}: at each of the $T$ occurrence events
    (processed in chronological order), the entity's full biography instance
    is rendered (\S\ref{sec:entity}) with the contested relation resolved to
    whichever side that event belongs to. Every background relation is
    rendered fresh at that same event regardless of which side fired -- this
    is what gives background content its exposure count of exactly $T$, with
    no separate scheduling required.
  \item \textbf{Logging}: the exact position and side-label of every
    occurrence event is recorded per entity, sufficient to reconstruct
    $(n_A, n_B)$ and total exposure count at analysis time.
\end{itemize}

\subsection{Corpus Scale: Conflict-Pair Phase}
\label{sec:scale}

\begin{itemize}
  \item \textbf{Split levels: $T=80$, 6 levels}: a working estimate reasoned
    from the capacity-slack argument of \S\ref{sec:seeds}, not itself a
    pilot-validated value. Allen-Zhu's $\sim$1000-exposure regime is the
    wrong anchor here (\S\ref{sec:seeds} shows this population is nowhere
    near capacity-constrained), but very low exposure counts likely
    undershoot the point where a contested pair's confidence gap is reliably
    learned rather than noisy. Split levels, all summing to $T=80$:
    $(40,40)$, $(50,30)$, $(60,20)$, $(68,12)$, $(73,7)$, $(77,3)$.
  \item \textbf{Replication}: 20 replicate entities per (relation type, split
    level) combination (14 types $\times$ 6 levels $\times$ 20 $=$ 1{,}680
    entities total). Relation type and split level are both exactly balanced
    by the construction in \S\ref{sec:entity}, so relation-type identity
    cannot confound the frequency-gap comparison; background-relation-type
    usage is balanced via the same cyclic-offset rule.
\end{itemize}

\subsection{Corpus Assembly}
\label{sec:assembly}

Interleaves the OpenWebText backbone (\S\ref{sec:corpus}) with entities'
vignette occurrences (\S\ref{sec:scheduler}) into the packed token shards
\texttt{training/data.py:PackedTokenDataset} expects (\texttt{meta.json} +
\texttt{*.bin}), plus a separate occurrence log.

\begin{itemize}
  \item \textbf{Single streaming pass}: each occurrence event's position
    (\S\ref{sec:scheduler}) is a real-valued coordinate into
    $[0, \texttt{total\_tokens})$, decided purely from the target token
    budget with no dependency on backbone content, so all events are
    precomputed and sorted up front, then fired online while streaming.
    \textbf{Two counters are kept, deliberately not one}: a backbone-only
    cumulative count (compared against event positions, which are defined
    against the backbone stream alone,
    \S\ref{sec:corpus-vs-training-tokens}) and the actual full
    packed-stream offset including already-injected tokens (used only for
    the occurrence log). Comparing events against the full-stream offset
    instead would let an injected vignette's own tokens push that counter
    past later events' positions, firing them in a cascade before the
    backbone content meant to precede them had been written.
  \item \textbf{Document boundaries and insertion granularity}: each
    OpenWebText row is one self-contained article and each injected
    vignette occurrence is its own self-contained document, both delimited
    by the GPT-2 EOS token. Injections are not restricted to firing only
    between whole articles, since OpenWebText documents run up to
    $\sim$100k characters -- checking only at document boundaries would let
    a scheduled position sit unfired for an entire long article. Instead
    each document is tokenized once at full length
    (\texttt{backbone\_openwebtext.py:tokenize\_with\_unit\_boundaries},
    which also returns paragraph- or sentence-boundary token indices), and
    the event-firing check runs after every such boundary. EOS is appended
    only after a document's final boundary, so an uninterrupted article
    stays one EOS-delimited unit; a vignette firing at an internal paragraph
    boundary gets an EOS forced immediately before it as well as after, so
    it never runs directly into the preceding backbone paragraph with no
    separator.
  \item \textbf{\texttt{total\_tokens} is a backbone budget, not the final
    corpus length}: backbone lines are tokenized and accumulated until this
    budget is met; scheduler positions are coordinates into that
    backbone-length stream. Injected occurrences add their own tokens on
    top, so the realized corpus is slightly longer than
    \texttt{total\_tokens} by the total injected-token count. This avoids a
    circular dependency, since injected-text length varies per
    LLM-authored variant and is not known until generation time.
  \item \label{sec:corpus-vs-training-tokens} \textbf{\texttt{total\_tokens}
    (on-disk corpus size) and total training tokens (\S\ref{sec:optimization})
    are independent quantities}: this section's \texttt{total\_tokens} fixes
    how many distinct tokens exist in the packed shards on disk. Total
    training tokens is \texttt{batch\_size} $\times$ \texttt{n\_positions}
    $\times$ \texttt{max\_steps} (\S\ref{sec:optimization}) -- the number of
    tokens the optimizer actually processes, which routinely exceeds
    \texttt{total\_tokens} because
    \texttt{training/data.py:PackedTokenDataset} is sampled with
    replacement: \texttt{training/train.py:infinite\_loader} reshuffles and
    re-cycles the same fixed set of on-disk windows indefinitely rather than
    stopping when the corpus runs out. A corpus with $N$ on-disk tokens
    trained for $M$ total training tokens revisits that corpus
    $\approx M/N$ times over the run; this ratio must stay small enough that
    repetition itself doesn't degrade quality
    (\citet{muennighoff2023scaling}: negligible loss up to $\sim$4 epochs, a
    $\sim$16-epoch half-life past which returns diminish sharply).
  \item \textbf{Chunked consumption}: each $\sim$2GB text chunk from
    \texttt{backbone\_openwebtext.py} is tokenized and its documents written
    directly to the output shard as they're produced, then discarded before
    the next chunk is pulled -- never holding more than one chunk's raw
    text, or the full token stream, in memory at once.
  \item \textbf{Injected-token floor}: the number of occurrence events is
    fixed by population size and $T$ (1{,}680 entities $\times$ $T$,
    independent of \texttt{total\_tokens} -- it only controls how spread out
    they are, not how many there are), so there is a fixed injected-token
    floor that \texttt{total\_tokens} must be large relative to for the
    ``rare conflicting fact in an otherwise natural corpus'' framing to
    hold. At $T=80$ (\S\ref{sec:scale}) and $\sim$103 tokens/occurrence
    (gpt-4o-mini vignettes): $1{,}680 \times 80 = 134{,}400$ occurrences,
    $\approx$13.8M injected tokens -- $\approx$0.55\% of the fixed
    $2.5\times10^9$ token budget (\S\ref{sec:corpus}).
  \item \textbf{Variant selection}: round-robin per (entity, side), same
    discipline as \S\ref{sec:entity} (\texttt{PoolCycler}) -- no variant
    repeats until that side's bank (\texttt{v\_a} or \texttt{v\_b}
    variants, 8 total per entity split unevenly per \S\ref{sec:entity}) is
    exhausted, then reshuffled. Heavy repetition is expected and accepted:
    at the high-occurrence side of an extreme split (e.g.\ $(77,3)$'s
    77-occurrence side with 7 variants), each paragraph is reused
    $\approx$11 times over that entity's exposure history. This is an
    explicit cost/API-budget tradeoff against \citet{allenzhu2024capacity}'s
    bioR anchor of 40 variants/entity.
  \item \textbf{Output}: the packed stream (\texttt{data/processed/<name>/})
    plus \texttt{results/occurrence\_log.json}, recording every occurrence
    event's entity, side, and final absolute token position in the packed
    stream -- sufficient to reconstruct $(n_A, n_B)$ and exposure timing at
    analysis time, and for the training loop to derive dense-checkpoint step
    windows (\S\ref{sec:checkpointing}) from token position and batch/block
    size.
  \item \textbf{Eval decoupling}: per-fact recall is \emph{not} scored
    against training vignettes. Side $A$ and side $B$ training text are
    independently LLM-authored and share no common stem, so there is no
    single shared token position to compare at. Instead, the
    per-relation-type template banks generated in \S\ref{sec:generation} are
    repurposed as a fixed, controlled, eval-only probing format: side $A$
    and side $B$ are rendered through the identical template, diverging at
    exactly one token as designed in \S\ref{sec:divergence}. Training text
    stays free-form and diverse; evaluation runs on a separate, minimal,
    controlled artifact the model never saw during training.
\end{itemize}

\subsection{Optimization}
\label{sec:optimization}

\begin{itemize}
  \item Optimizer: AdamW (fused CUDA kernel when running on GPU), cosine LR
    schedule with warmup, weight decay 0.01.
  \item Precision: bf16 mixed precision, gradient clipping.
  \item \textbf{Batch size / step count: 128 / 38{,}146}. An empirical CUDA
    probe (\texttt{training/probe\_max\_batch\_size.py}) found the true max
    batch size that fits on an A100-80GB at 512-token context is 140;
    128 leaves $\approx$26GB of headroom, comfortably clear of the backward-pass
    logits buffer that dominates peak memory. Architecture and weight
    decay/cosine schedule shape match \citet{allenzhu2023physics} for
    comparability; batch size, grad\_clip, and AdamW eps do not (below). Step
    count is derived from the fixed $2.5\times10^9$-token target
    (\S\ref{sec:corpus}): at 512-token context, $128 \times 512 \times
    38{,}146 \approx 2.5\times10^9$ tokens, matching
    \texttt{total\_tokens} (\S\ref{sec:corpus-vs-training-tokens}) so this
    remains $\approx$one epoch over the backbone.
  \item \textbf{Learning rate and stability settings: lr $6\times10^{-4}$,
    warmup\_steps 2000, grad\_clip 0.5, eps $10^{-6}$}. lr follows the
    standard GPT-2-small/GPT-3-125M literature value
    \citep{radford2019language, brown2020gpt3}. warmup covers $\sim$5\% of
    \texttt{max\_steps}; grad\_clip is tightened and eps set explicitly (not
    AdamW's implicit $10^{-8}$ default) per \citet{wortsman2024smallscale}'s
    finding that a too-small eps can let a layer's effective LR blow up once
    gradient RMS shrinks under high-LR instability. These settings, together
    with the checkpoint-rotation/spike-guard/rollback mechanism of
    \S\ref{sec:checkpointing}, were adopted after early training runs at a
    more aggressive, sqrt-rescaled learning rate suffered irrecoverable loss
    spikes partway through training; the literature-standard lr combined
    with a tighter eps/grad\_clip has since trained cleanly across the full
    $T$-sweep (\S\ref{sec:seeds}).
  \item Estimated wall-clock time on a single A100 80GB: using nanoGPT's
    published single-A100 benchmark ($\sim$178,000 tokens/sec at 49\% MFU
    for this same 124M-parameter architecture, measured with a fused
    optimizer and \texttt{torch.compile}, both also used here) as the
    closest available reference point, $2.5\times10^9$ tokens $/ 178{,}000
    \approx 3.9$ hours -- comfortably inside a single 12-hour cluster job
    (CLAUDE.md \S2), treated as a planning estimate rather than a measured
    commitment.
\end{itemize}

\subsection{Checkpointing}
\label{sec:checkpointing}

Rotating slots, step-indexed, with an automatic stability layer that
supersedes the wall-clock checkpoint rule in CLAUDE.md \S5:
\begin{itemize}
  \item \textbf{Sparse}: full checkpoint (weights, optimizer state, RNG
    state) every $S$ steps, $S$ set to bound compute loss on crash to
    $\sim$30 minutes at measured throughput -- written to a free
    \texttt{slot\_N} (crash-safe temp-write-then-rename) until
    \texttt{num\_checkpoint\_slots}$=3$ exist, then evicts the
    oldest-by-step slot.
  \item \textbf{Dense}: every few hundred steps within a fixed window
    following each contested-relation occurrence event
    (\S\ref{sec:scheduler}), logs eval metrics only; no full weights are
    captured on this cadence.
  \item \textbf{Spike guard}: per training step, a 50-step rolling mean of
    loss is tracked; if a step's loss is non-finite or exceeds $4\times$
    that mean (once at least 20 steps of history exist), the optimizer
    update is skipped for that step (gradients zeroed rather than applied)
    and the spiked loss is excluded from the rolling mean.
  \item \textbf{Rollback}: after each eval, a streak counter increments
    whenever val\_ppl exceeds $2\times$ the best val\_ppl seen so far; once
    that streak reaches 3 consecutive evals, model/optimizer/scheduler/RNG
    state is reloaded from the most recent rotating checkpoint, the
    learning rate is multiplied by a $0.7\times$ decay factor (floor: 0.1 of
    base lr), and the step counter is rewound to match. Rollback is capped
    at 5 occurrences; once exhausted, training stops early rather than
    looping indefinitely, and the run summary records that it ended via
    exhausted-rollback rather than by reaching \texttt{max\_steps}. This
    state is process-local: it does not persist across a separate
    preemption/resume of the training script.
  \item \textbf{Final}: an unconditional full save at the true last training
    step, through the same rotating-slot mechanism as Sparse above.
  \item \textbf{Finalize}: once training ends,
    \texttt{training/finalize\_checkpoint.py} renames the single
    lowest-val\_ppl retained slot to \texttt{checkpoints/best/} and deletes
    every other slot, so exactly one checkpoint remains on disk per run,
    and that is what any downstream eval (\S\ref{sec:eval}) reads by
    default.
\end{itemize}

\subsection{Evaluation}
\label{sec:eval}

Run at every checkpoint:
\begin{itemize}
  \item Held-out perplexity against an in-domain OpenWebText validation
    split (LM coherence check, used for checkpoint selection); held-out
    WikiText-103 perplexity is additionally reported as an out-of-domain
    diagnostic.
  \item Per-fact recall: log-probability / rank of the target value's first
    divergent token (\S\ref{sec:divergence}), computed per side ($A$/$B$) of
    the contested relation for every entity. Scored at a single token
    rather than the full continuation to avoid a token-count length
    confound, and via log-probability rather than free-form generation
    because base (non-instruction-tuned) checkpoints cannot be reliably
    prompted for parseable short-answer output.
\end{itemize}

\textbf{Non-conflict (background) recall test} (\texttt{eval/recall.py}): a
dedicated evaluation of the 4 background (non-contested) facts per entity,
separate from the contested-term calibration question above -- background
facts receive exactly $T$ exposures per entity regardless of contested split
level (\S\ref{sec:entity}), giving one clean recall-vs-$T$ curve per relation
type, and serving as a check that training produced fact learning at all,
independent of the harder frequency-calibration question. Uses the same
fixed-template probing stem as contested recall (\S\ref{sec:divergence}), but
scored against the \emph{whole} $50{,}257$-token vocabulary at absolute
probability, so the model must outscore every plausible completion, not just
other entries from that relation's value pool. Metrics, per relation type and
pooled overall across all $6{,}720$ background probes ($1{,}680$ entities
$\times$ 4 facts): \textbf{top-1 accuracy}, \textbf{top-5 accuracy}, and
\textbf{logit margin} (signed: the model's winning margin over the runner-up
when its top pick is correct, or the deficit behind its actual top pick when
wrong). The pooled figure is compared across the $T$-sweep
(\S\ref{sec:seeds}); per-record \texttt{relation\_key}/\texttt{entity\_id}
tags are retained in the output JSON so results remain sliceable by relation
type after the fact.

\subsection{Calibration Target}
\label{sec:calibration}

Both sides of an intra-memory conflicting pair are model-internal claims with
no external ground truth, so calibration cannot be assessed against ``which
answer is correct.'' The target is instead defined relative to the exposure
statistics the corpus generator itself controls (frequency,
\S\ref{sec:corpus}): a rational aggregator of that evidence should produce a
confidence over $\{A, B\}$ satisfying three checkable, qualitative properties
rather than an exact posited functional form (e.g.\ an exact Bayesian
posterior over exposure counts):

\begin{itemize}
  \item \textbf{Monotonicity}: confidence in the higher-frequency side moves
    monotonically as the frequency gap $n_A - n_B$ widens; no reversals as
    the gap is dialed up.
  \item \textbf{Symmetry at balance}: at the balanced condition ($n_A = n_B$)
    confidence sits near 50/50 rather than locking arbitrarily onto one
    side.
  \item \textbf{No premature saturation}: confidence does not approach the
    softmax extremes long before the frequency gap does (overconfidence).
\end{itemize}

\textbf{Metrics}, derived from the same per-side log-probabilities already
scored in \S\ref{sec:eval} (no separate eval pipeline): softmax the two
log-probabilities into a confidence over $\{A, B\}$; plot against the
frequency gap over the split-level sweep of \S\ref{sec:scale}; compute
expected calibration error against the chosen normative curve, and a
monotonicity-violation count. By default this is computed from the
\emph{logit gap} -- the raw logit difference between the two sides at their
divergent token (\S\ref{sec:divergence}), equivalent up to a sigmoid
transform to the softmax confidence above -- and both the gap and the full
per-side logits are recorded, so alternative derived metrics remain
reconstructable without rerunning the model.

This metric is scored at the base (non-instruction-tuned) checkpoint stage
per \S\ref{sec:eval}, which keeps the token-level probability distribution
itself informative -- the object this section's metrics operate on -- in
addition to being a practical necessity, since base checkpoints cannot be
reliably prompted for parseable short-answer output.

\subsection{Candidate Quantitative Metric: Cross-Entropy to a Proportional Target}
\label{sec:xent-metric}

The three properties above are checkable but not optimizable: they describe
what a well-calibrated confidence curve should look like, not how far a
given curve is from that in a single number. This section defines one
candidate quantitative operationalization, among several under active
consideration; a Bayesian posterior-mean target was also tried and is
provisionally ruled out below.

\textbf{Target}: for a contested entity with exposure split $(n_A, n_B)$,
$T = n_A + n_B$, define the ideal post-softmax distribution over $\{A, B\}$
as directly proportional to the exposure split: $q_A = n_A / T$,
$q_B = n_B / T$ -- the maximum-likelihood target under the simplest possible
reading of the corpus as evidence (each occurrence is one vote for its side,
no prior pull toward uncertainty).

\textbf{Metric}: let $p_A = \sigma(\ell_A - \ell_B)$ be the model's actual
softmax confidence restricted to the two candidates (the logit gap above,
sigmoid-transformed), $p_B = 1 - p_A$. Score the model's deviation from the
target via cross-entropy, $H(q, p) = -q_A \log p_A - q_B \log p_B$, reported
per split level (mean over entities at that split) and pooled. The related KL
divergence $D_{KL}(q \| p) = H(q,p) - H(q)$, where
$H(q) = -q_A\log q_A - q_B\log q_B$ is the target's own entropy, is computed
at negligible extra cost and reported alongside it: $H(q)$ varies by split
level (maximal at balance, near zero at the most extreme split), so raw
cross-entropy conflates irreducible target uncertainty with actual
miscalibration, while $D_{KL}$ isolates the latter. Cross-entropy is reported
as the primary number, $D_{KL}$ as the split-level-comparable variant.

\textbf{Caveat}: both quantities are computed per entity then averaged, which
conflates whether the \emph{average} confidence matches the average target
(systematic bias) with how much individual-entity confidence scatters around
that average (per-entity noise). By Jensen's inequality (KL is convex),
scatter alone inflates the pooled number even when the aggregate is
well-calibrated; empirically, at the balanced split, pooled mean confidence
(53.4\%) sits close to the 50\% target but mean KL is not the smallest of the
six split levels. Disentangling the two is left as future work.

\textbf{Implementation}: computed post-hoc from already-recorded per-entity
$\ell_A, \ell_B$ (no rerun of the model needed) in \texttt{eval/recall.py}'s
\texttt{\_summarize\_contested}, reported for the relation-type-restricted
population of \S\ref{sec:relation-restriction} by
\texttt{eval/build\_intervention\_baseline.py} under
\texttt{results/interventions/<run-name>/}.

\textbf{Example numbers} (T=1280, top-7 relations, checkpoint step 61035;
\texttt{results/interventions/gpt2-small-openwebtext-T1280/
calibration\_top7\_step61035.md}): pooled mean cross-entropy 0.875 nats,
pooled mean $D_{KL}$ 0.409 nats. By split level ($n_A$/$n_B$, mean
cross-entropy, mean $D_{KL}$): 640/640: 1.255, 0.562; 800/480: 1.204, 0.542;
960/320: 1.115, 0.553; 1088/192: 0.849, 0.427; 1168/112: 0.551, 0.255;
1232/48: 0.275, 0.115 -- both quantities fall roughly monotonically as the
split widens, even though the raw logit gap \emph{overshoots} the
proportional target at the two widest splits (an informal fit-comparison
against $\log(n_A/n_B)$ found the same overshoot and ruled out Bayesian
posterior-mean shrinkage as an explanation, since any valid prior can only
pull the target toward 50/50, the wrong direction to explain an overshoot).
Cross-entropy/KL treat ``more confident than the target, on the correct
side'' as cheap once $q$ is already near an extreme, unlike a raw-logit-gap
comparison -- exactly the kind of behavior difference between candidate
metrics this exploration phase is meant to surface.

\subsection{Exposure-Count Calibration and Seeds}
\label{sec:seeds}

\textbf{Exposure-budget pilot}: rather than adopting an exposure-count regime
from prior work, $T$ (\S\ref{sec:scale}) is meant to be calibrated
empirically to this project's own setting. \citet{allenzhu2024capacity}'s
$\sim$1000-exposure regime is a poor prior here: their numbers are
calibrated to a capacity-stressed setting (millions of entities competing
for a fixed parameter budget), whereas this design uses a much smaller,
non-capacity-constrained entity population, so any single fact is learned
far faster and values near 1000 exposures would be expected to saturate both
sides of a contested pair well before the frequency gap does -- directly
violating the ``no premature saturation'' property of \S\ref{sec:calibration}.
The pilot itself has not been run yet; in its place, $T=80$
(\S\ref{sec:scale}) is a working estimate reasoned from a capacity argument:
this population's total factual content is
$1{,}680 \times 5 \times \log_2(150) \approx 61{,}000$ bits (5 facts/entity,
$\sim$150-entry value pools), against GPT-2-small's $\sim$2-bit/parameter
capacity ceiling of $124\text{M} \times 2 \approx 248$M bits
(\citet{allenzhu2024capacity}'s own measured saturation rate) -- a
$\sim$4{,}000$\times$ slack, i.e.\ nowhere near capacity-constrained, so
substantially fewer than 1000 exposures should suffice. Once the pilot is
run, it should sweep a range spanning this estimate (e.g.\
$32/64/128/256/512$) and inspect the resulting logit-gap/calibration curve;
$T$ is fixed at a value where that curve is monotonic but not yet saturated,
per \S\ref{sec:calibration}. Only the conditions selected by this pilot are
then re-run at the full corpus scale and seed count of \S\ref{sec:scale}.

\textbf{Seeds}: multiple random seeds and data-ordering variants per
condition, to separate a stable property of the setup from a
single-trajectory artifact, staged the same way as the exposure-budget pilot
above (single-seed pilot identifies conditions of interest; only those are
re-run at full seed count).

\textbf{Pilot scope}: of the original open parameters, $T$ remains the one
genuinely unresolved:
\begin{itemize}
  \item \textbf{\texttt{total\_tokens}} (\S\ref{sec:corpus},
    \S\ref{sec:assembly}): fixed at $2.5\times10^9$, the Chinchilla-optimal
    point for a 124M-parameter model \citep{hoffmann2022chinchilla}, now
    that OpenWebText (\S\ref{sec:corpus}) can supply that many unique tokens
    in a single pass.
  \item \textbf{Backbone exposure}: with \texttt{total\_tokens} set equal to
    the total-training-token target, the backbone is seen $\approx$once (no
    repetition), resolved by construction rather than by a separate ratio
    decision.
  \item \textbf{$V$} (variants per entity/side, \S\ref{sec:entity}): fixed
    at 8 total per entity, split unevenly across sides via
    \texttt{variants\_per\_side} (\S\ref{sec:entity}, \S\ref{sec:scale}),
    still well short of \citet{allenzhu2024capacity}'s bioR anchor
    (40/entity); heavy round-robin repetition on the high-occurrence side of
    extreme splits is accepted as a cost/API-budget tradeoff
    (\S\ref{sec:assembly}).
  \item \textbf{$T$} (\S\ref{sec:scale}): $80$ is a reasoned working
    estimate (capacity-slack argument above), not a pilot-validated value.
    The pilot, once run, should sweep around this estimate and may revise
    $T$ (and, if $T$ moves far enough, revisit $V$'s allocation and the
    split levels of \S\ref{sec:scale}).
\end{itemize}

\subsection{Recall Baseline and Relation-Type Restriction for Interventions}
\label{sec:relation-restriction}

Every method surveyed below arbitrates between two \emph{already-memorized}
candidate values for a contested relation -- steering, pruning, or
re-calibrating confidence over $\{A, B\}$ presupposes that both sides were
actually stored during training, not merely that the eval template happens
to match training phrasing. Until this is checked per relation type, none of
these methods has a valid baseline to act on: an intervention ``working'' on
a relation the model never learned the fact for is indistinguishable from an
intervention doing nothing to two tokens the model was already guessing at
chance.

The background-recall diagnostic (\S\ref{sec:eval}), rerun with
phrasing-tolerant best-of-5-template scoring on the $T{=}1280$ checkpoint
(\S\ref{sec:seeds}) to rule out template mismatch as a confound, shows
recall is not uniform across the 14-type inventory (\S\ref{sec:entity}) but
sharply bimodal:

\begin{itemize}
  \item \textbf{Retained (top 7, best-of-5 top-1 $\geq 65\%$)}: birthplace
    ($99.2\%$), field of expertise ($99.4\%$), current residence ($93.8\%$),
    employer/role ($89.6\%$), alma mater ($77.7\%$), mentor ($70.6\%$),
    license/certification ($65.6\%$).
  \item \textbf{Excluded (bottom 7, best-of-5 top-1 $\leq 43.3\%$)}: working
    language ($43.3\%$), publication venue ($23.5\%$), civic role ($8.1\%$),
    authored work ($7.3\%$), affiliation ($0.8\%$), funding source
    ($0.2\%$), award or honor ($0.0\%$).
\end{itemize}

Five of the bottom seven stay under $8\%$ even under best-of-5 scoring,
which was specifically designed to absorb template/training-phrasing
mismatch -- so this is not a scoring artifact but a genuine failure to store
(or extract) the fact, at the same $T{=}1280$ exposure count ($1.28\times$
\citet{allenzhu2024capacity}'s own 1000-exposure calibration point) that
yields near-ceiling recall for the top-7 group. Capacity is therefore not
the binding constraint (contra the slack argument of \S\ref{sec:seeds}) --
this instead matches \citet{allenzhu2024capacity}'s storage-vs-extraction
distinction (their Part 3.1): the same model, same capacity budget, same
exposure count, stores some relation types reliably and others not at all.
Root-causing \emph{why} these particular five relation types fail
(value-string structure was checked directly and does not distinguish the
two groups) is left as future work, not a blocker for intervention
experiments on the top 7.

\textbf{Scope of the restriction}: intervention experiments below restrict
the contested-relation-type inventory (\S\ref{sec:entity}) to the top 7
only -- entities whose contested relation is one of the bottom 7 are dropped
from the intervention population. This is a post-hoc filter applied to the
existing $T{=}1280$ corpus and checkpoint, not a corpus regeneration:
background relations are unaffected (reinforcement only, not an analyzed
variable, \S\ref{sec:entity}), and no retraining is required.

\section{Baseline Methods}

Baselines are grouped by mechanism. This project's setting is
\emph{intra-memory} conflict: both candidate answers are already parametric,
so there is no external context to contrast against a prior. Methods
originally designed for \emph{context-memory} conflict (context vs.\
parametric) are marked as such; applying them here requires substituting a
disambiguating cue (e.g., a time qualifier) for the missing external context.
Every method below operates on the relation-type-restricted entity
population of \S\ref{sec:relation-restriction} -- the top 7 relation types
with confirmed reliable recall, not the full 14-type inventory.

\subsection{Logit/Decoding-Level Contrastive Methods}
\label{sec:decoding}

Output-distribution-only methods, no internal access required. All but DoLa
presuppose a with-context / without-context contrast, substituted here with a
disambiguating-cue pair.

\begin{description}
  \item[CAD] \citep{shi2023trusting} Contrasts with- and without-context
    output distributions with a fixed amplification coefficient.
  \item[COIECD] \citep{yuan2024discerning} Token-level entropy gate on CAD;
    contrastive adjustment applied only where an entropy-based conflict
    signal fires.
  \item[AdaCAD] \citep{wang2024adacad} Replaces CAD's fixed coefficient with a
    per-step, Jensen--Shannon-divergence-derived weight.
  \item[CoCoA] \citep{khandelwal2025cocoa} Adds confidence-aware terms
    (entropy gap, contextual peakedness) to AdaCAD's divergence-based weight.
  \item[DoLa] \citep{chuang2023dola} Contrasts early/middle-layer and
    final-layer logit distributions within a single forward pass; no context
    required. Directly applicable to the intra-memory setting without
    adaptation.
  \item[ARR] \citep{jiang2026contextaware} Routes each decoding step between
    trust-prior, trust-context, and agree regimes based on the conflict
    signal, rather than presupposing context trustworthiness as CAD/COIECD/
    AdaCAD/CoCoA do. Closest existing precedent for arbitrating between two
    parametric candidates without an external anchor.
\end{description}

\subsection{Activation Steering Methods}
\label{sec:steering}

Vector-based residual-stream interventions. Contrastive-pair construction does
not require external context, so this family ports to the intra-memory
setting more directly than \S\ref{sec:decoding}: pairs are built from two
disambiguated prompts, each eliciting one side of a conflicting-fact pair.

\begin{description}
  \item[CAA] \citep{rimsky2024steering} Steering vector from the
    mean-activation difference between contrastive prompt pairs, added
    (scaled) into the residual stream. Reference baseline for this family.
  \item[SpARE] \citep{zhao2024spare} Steers sparse-autoencoder features rather
    than dense residual directions; reports the conflict signal is linearly
    decodable from mid-layer activations onward, in the context-memory
    setting. Decodability under intra-memory conflict (no context-derived
    feature to isolate) is untested and should be treated as a hypothesis.
  \item[K-CAST] \citep{valentino2025mitigating} kNN-based conditional
    steering: per-instance direction from nearest-neighbor lookup rather than
    a single fixed vector. Originally applied to a reasoning/content-effects
    task.
  \item[ContextFocus] \citep{anand2026contextfocus} Context-faithfulness
    steering; reports composability with prompting-based methods. Requires
    the most adaptation of this group, being a context-faithfulness method by
    construction.
\end{description}

\subsection{Architecture/Attention-Level Intervention}
\label{sec:attention}

\begin{description}
  \item[PH3] \citep{jin2024cutting} Path-patching identification of attention
    heads causally responsible for propagating a memorized fact, followed by
    pruning/down-weighting at inference time. No context-vs-parametric
    contrast required; with DoLa, one of two methods surveyed here that is
    intra-memory-native by construction. Closest prior art to this project's
    CMAP intervention.
\end{description}

\subsection{Conflict-Detection Signals}
\label{sec:detection}

Detection only -- no output-changing mechanism.

\begin{description}
  \item[DynamicQA] \citep{marjanovic2024dynamicqa} Semantic entropy over
    paraphrase completions (adapted from Kuhn et al.) as an intra-memory
    conflict estimator; reports higher conflict for temporally/disputably
    dynamic facts than for static facts.
  \item[Self-detection] \citep{zhao2024knowing} Answer divergence across
    paraphrased questions (``memory strength'') as a signal for questions the
    model does not reliably know.\footnote{Yukun Zhao et al.; not to be
    confused with Yu Zhao, first author of SpARE (\S\ref{sec:steering}).}
\end{description}
Either signal is a candidate gate for the causal interventions in
\S\ref{sec:attention}, e.g.\ restricting PH3/CMAP-style pruning to
high-entropy tokens. The calibration target of \S\ref{sec:calibration}
additionally licenses a comparison unavailable to either source paper, which
lacks access to ground-truth exposure statistics: whether semantic entropy or
answer-divergence tracks the true frequency evidence gap as well as, or
better than, raw token-level confidence does.

\subsection{Calibration-Targeted Interventions}
\label{sec:calibration-interventions}

Methods that act on the calibration target of \S\ref{sec:calibration}
directly -- shaping how confident the model is -- rather than on which side
wins.

\begin{description}
  \item[Confidence penalty] \citep{pereyra2017regularizing} Training-time loss
    term penalizing low-entropy (over-peaked) output distributions, applied at
    the token positions where the two sides of a conflicting pair diverge. No
    architecture change, no context/cue required.
  \item[Temperature scaling] \citep{guo2017calibration} Post-hoc, single-scalar
    logit rescaling fit after training. Inference-time only, does not alter
    the model's factual content; the natural minimal baseline against which
    confidence-penalty training is compared.
  \item[Minimal in-context cue] Restating one side of the conflicting pair
    immediately before the query -- scaled down from the disambiguating-cue
    baselines of \S\ref{sec:decoding} to a form a base (non-instruction-tuned)
    model can plausibly use via induction-head-like copying rather than
    instruction comprehension. Tests whether a single in-context repetition
    shifts confidence consistently with, or in override of, the training-time
    exposure signal.
\end{description}

\subsection{Sampling/Ensemble-Based Mitigation}
\label{sec:sampling}

\begin{description}
  \item[Self-consistency] \citep{wang2022selfconsistency} Majority vote over
    multiple sampled completions; originally for chain-of-thought reasoning.
    Applied to factual recall: sample paraphrases of a probe question, vote
    or abstain on disagreement. Not evaluated in the literature as an
    intra-memory mitigation method specifically (only as a detection signal,
    \S\ref{sec:detection}); requires multiple forward passes per query, no
    internal access.
\end{description}

\section{Mechanistic Interpretability: Localization and Causal Control}
\label{sec:mechinterp}

\S\ref{sec:decoding}--\S\ref{sec:sampling} ask whether a given method shifts
the model's output toward one side of a contested pair; this section asks a
different, prior question -- \emph{how} the model stores and arbitrates
between the two sides internally -- and targets a causal, continuously
controllable intervention rather than a binary flip. Status: a rough plan
under consideration, not yet run or resourced; recorded here per CLAUDE.md
\S5 before any job is scoped or submitted.

\textbf{Relation to Baseline Methods}: this is a new analysis axis, not a
ninth baseline. It shares the relation-type-restricted population and
$T{=}1280$ checkpoint of \S\ref{sec:relation-restriction} (both sides of a
contested pair must already be confirmed stored for a localization question
to be meaningful), and reuses the disambiguated-contrastive-pair construction
of \S\ref{sec:steering} for cache/probe inputs. Everything downstream of that
-- layerwise probing, causal tracing, path patching, and dial-style steering
built from a causally-\emph{identified} locus rather than a fixed a-priori
layer -- is new. Implementation lives under \texttt{mech\_interp/}, kept
separate from \texttt{inference\_time/} (\S\ref{sec:decoding}--\S\ref{sec:sampling},
which already exists and ships candidate mitigations) and
\texttt{training\_time/}: \texttt{mech\_interp/} code reads a fixed, frozen
checkpoint and writes to \texttt{results/mech\_interp/}, and does not itself
select or ship a mitigation.

\subsection{Phase 0: Behavioral Scaffolding}
\label{sec:mechinterp-scaffolding}

Cheap, run first, and only correlational-vs-causal expectation-setting -- no
activation access required.

\begin{itemize}
  \item For each contested pair, build a generation probe (matching
    \S\ref{sec:eval}'s existing per-fact recall scoring) and a verification
    probe (forced-choice true/false over each side) for both $A$ and $B$.
  \item \textbf{Multi-verse check}: does the model \emph{generate} one side
    but \emph{verify} both as true? \citet{davidson2026futureoffacts} reports
    verification is learned before generation and survives continual
    learning more robustly, producing exactly this pattern -- a superseded
    side stays verified true even once it is no longer generated. Finding
    this pattern here (vs.\ clean suppression, where the losing side is
    verified false too) sets the baseline expectation for every later phase:
    a multi-verse result predicts Phase 3 will find full-fidelity
    dual-storage, while clean suppression is more consistent with one side
    being genuinely overwritten or never separately represented.
\end{itemize}

\subsection{Phase 1: Representational Localization (Correlational)}
\label{sec:mechinterp-reploc}

\begin{itemize}
  \item At every layer and at the key token positions (entity-mention token,
    final token before the divergent-token scoring position of
    \S\ref{sec:divergence}), train linear probes separately for ``does this
    activation encode $A$'' and ``does this activation encode $B$,'' using
    the contrastive-pair construction of \S\ref{sec:steering}.
  \item Compute the angle between the two probes' weight directions at each
    layer/position. Near-orthogonal is evidence for separable subspaces;
    small angle / high cosine overlap is evidence for superposition.
  \item Also probe for (i) a recency/exposure-order direction, per
    \citet{krasheninnikov2025freshinmemory}'s finding that training-order
    recency is linearly decodable from activations, and (ii) a
    conflict-ratio direction, keyed to the exposure split $(n_A, n_B)$ of
    \S\ref{sec:scale} and the proportional target $q_A = n_A/T$ of
    \S\ref{sec:xent-metric}. Check whether ``which side wins'' correlates
    with either of these directions rather than with a dedicated
    fact-identity direction -- a way this phase's localization claim could
    turn out to be confounded by a simpler, already-characterized signal.
\end{itemize}

\subsection{Phase 2: Causal Localization}
\label{sec:mechinterp-causal}

The causal test proper, following the activation-patching / causal-tracing
paradigm of \citet{meng2022rome} (localizing a single factual association)
and \citet{meng2023memit} (scaling the same causal signal to many facts):
run the model on an ``$A$'' context, cache all layer activations; run on the
matched ``$B$'' context; patch $A$'s cached activations into the $B$ run (and
the reverse), one layer/position at a time, measuring how much the output
flips back.

\begin{itemize}
  \item \textbf{Three granularities}: (i) full residual stream at a layer;
    (ii) individual attention heads, via path patching -- the same technique
    \S\ref{sec:attention}'s PH3 baseline \citep{jin2024cutting} already
    applies to knowledge-conflict head attribution, repurposed here for
    localization rather than pruning; (iii) individual MLP
    neurons/blocks, per ROME/MEMIT's own localization result that factual
    association storage concentrates in mid-layer MLPs.
  \item \textbf{Output}: a causal map of which components are
    \emph{necessary} (ablating kills the effect) and \emph{sufficient}
    (patching alone reproduces it) for expressing $A$ over $B$.
  \item \textbf{Superposition-vs-separation test}: if the same small
    head/neuron set drives both $A$-selection and $B$-selection (opposite
    sign or direction on the same component), that is superposition in one
    component. If disjoint component sets each causally control one side
    independently, that is closer to separate storage. This is the direct
    causal counterpart to Phase 1's angle measurement.
\end{itemize}

\subsection{Phase 3: Suppression vs. Erasure}
\label{sec:mechinterp-suppression}

Using the causal component(s) identified in Phase 2: ablate/patch them, then
apply the logit lens \citep{nostalgebraist2020logitlens} at earlier layers to
check whether the suppressed side is still fully recoverable elsewhere in the
network, as in ROME-edit's own ``suppression not erasure'' finding
\citep{meng2022rome}. This directly resolves the project's earlier
skepticism about ``both sides showing up in the logits'': confirming
recoverability upgrades that observation from \emph{both are in the logits}
to \emph{both are represented at full fidelity pre-suppression, and a
specific identified component performs the suppression}.

\subsection{Phase 4: Continuous Control}
\label{sec:mechinterp-dial}

Turns Phase 2's flip into a dial, and is the target deliverable of this
section (per \S\ref{sec:mechinterp}'s framing).

\begin{itemize}
  \item Build a steering vector from the difference between ``expressing
    $A$'' and ``expressing $B$'' activations, taken \emph{at the causal locus
    identified in Phase 2} -- unlike \S\ref{sec:steering}'s CAA baseline,
    which steers at a layer fixed a priori rather than one selected by a
    causal test.
  \item Scale the vector continuously; record the model's output
    distribution over $\{A, B\}$ as a function of steering magnitude.
  \item Compare the resulting curve against the true exposure-ratio target
    $q_A = n_A/T$ of \S\ref{sec:calibration}/\S\ref{sec:xent-metric}. A curve
    that can be dialed to hit any target ratio smoothly is evidence for a
    blended/superposed representation; discrete jumps are evidence for
    separately stored facts behind a threshold-like selector.
\end{itemize}

\subsection{Phase 5: Sweep Across Training-Time Variables}
\label{sec:mechinterp-sweep}

Repeat Phases 1--4 across \S\ref{sec:scale}'s split-level sweep, and, once
run, \S\ref{sec:seeds}'s exposure-budget ($T$) pilot, to test whether the
\emph{type} of mechanism (separate vs.\ superposed; sharp vs.\ continuous)
itself depends on these training statistics -- reconnecting the
training-dynamics thread of \S\ref{sec:scale}--\S\ref{sec:seeds} to the
localization thread of this section. This is the section's most likely
source of a genuinely novel claim if it holds, e.g.\ that conflicts
introduced with a wide early split resolve via separate localized
components while narrow late-forming splits resolve via
superposition/suppression.

\subsection{Controls and Reporting}
\label{sec:mechinterp-controls}

\begin{itemize}
  \item Matched clean/corrupt example pairs and random-head/neuron baselines
    throughout, so effect sizes are interpretable against chance rather than
    against an unnormalized raw patching effect.
  \item At GPT-2-small's scale, full path patching across all heads is
    tractable without the head-subsampling large-model mechanistic
    interpretability work typically requires.
  \item Report both the causal effect size, in the sense of
    \citet{elazar2022measuring}'s causal framework, and Phase 1's
    probe-orthogonality result side by side: they test the same
    separation-vs-superposition hypothesis from causal and correlational
    angles respectively, and agreement between them is the strongest
    available evidence for whichever answer they converge on.
\end{itemize}

\subsection{Open Scope Items}
\label{sec:mechinterp-module}

Not yet resourced or scheduled -- to be confirmed per CLAUDE.md \S5 before
any job is submitted:
\begin{itemize}
  \item Which relation types within the top-7 restriction
    (\S\ref{sec:relation-restriction}) to prioritize for the initial pass,
    vs.\ running all seven from the start.
  \item Layer/position sampling budget for Phase 1's probe sweep and Phase
    2's path-patching sweep (both tractable in full at this model scale per
    \S\ref{sec:mechinterp-controls}, but full-grid runtime has not been
    measured).
  \item Whether Phase 5's sweep runs first at $T{=}1280$ alone, or extends
    to the $T{=}80$/$T{=}320$ checkpoints already on disk
    (\S\ref{sec:seeds}) in the same pass.
\end{itemize}

\bibliographystyle{plainnat}
\bibliography{references}

\end{document}
